\section{Background and Related Work}
\label{sec:related-work}

The literature reviewed below was identified through a structured, non-systematic search of gallbladder-ultrasound AI and medical-XAI publications (2016--2026). It followed the spirit of search-transparency practice recommended for systematic reviews \cite{ref15}. This is a narrative review, not a formal systematic one.

Conventional CAD systems for gallbladder imaging relied on hand-crafted texture features and classical classifiers \cite{ref16}. Deep learning has since substantially improved performance on this task. Turki et~al. \cite{ref17} released the multi-class UIdataGB gallbladder ultrasound dataset used in this paper. It spans nine disease classes drawn in part from ultrasound video sequences. Their dataset-release publication reports that the underlying images were collected across three to four hospitals in Baghdad, Iraq, over four years, on four different ultrasound scanner models. The public release itself, still, carries no hospital, patient, or scanner identifier. So this multi-site provenance cannot be recovered or stratified on from the released files (Section~\ref{sec:methods}.1).

Using the same 10,692-image, 1,782-patient collection, Obaid et~al. \cite{ref18} reported 98.35\% accuracy with a MobileNet classifier across all nine classes. They evaluated it with a single CAM variant, reporting the results qualitatively rather than through a quantitative, multi-method XAI comparison. Two more recent studies on this dataset use more than one XAI method. Nabil et~al.'s MSFE-GallNet-X \cite{ref57} applies both Grad-CAM and LIME. Kumar et~al. \cite{ref58}, on a related cholelithiasis-focused ultrasound task, combine Grad-CAM and LIME into a hybrid explanation. Both report their explanations qualitatively, as illustrative heatmaps, rather than with the quantitative cross-method, cross-replicate, or faithfulness metrics applied here. They show that multiple XAI methods have been used together for gallbladder ultrasound. This narrows, but does not close, the gap this paper addresses (see the gap statement below).

A parallel line of work has built classifiers on a separate, five-class gallbladder-ultrasound dataset (GBCU). Basu et~al. \cite{ref7} released GBCU alongside curriculum-learning ResNet classification, and RadFormer \cite{ref19} added global-local transformer attention for interpretable GBC detection on the same dataset. Gupta et~al. prospectively evaluated deep learning for GBC detection on ultrasound \cite{ref20}. They also studied a multiple-instance-learning approach across several explicitly identified centres \cite{ref21}. This kind of externally validated, site-stratified evaluation is something this paper's single, site-unlabelled dataset release cannot support. More recent 2024--2025 work has pushed accuracy further. Anatomy-aware architectures on GBCU appear in Hasan et~al.'s GBCHV \cite{ref22}. Video-based masked autoencoder detection shows up in Basu et~al.'s FocusMAE \cite{ref23}, building on an earlier weakly-supervised variant \cite{ref24}. Other work addresses gallbladder polyp classification specifically \cite{ref16}, \cite{ref25}, \cite{ref26}. With the partial exception of RadFormer's attention maps, none of these studies, on either dataset, compares multiple XAI methods quantitatively or tests reproducibility across independent training runs.

Deep convolutional networks are now the dominant approach across medical-imaging classification generally. And transfer learning from ImageNet-pretrained backbones such as ResNet \cite{ref27}, DenseNet \cite{ref28}, and EfficientNet \cite{ref29} is standard practice for small medical datasets \cite{ref30}. That performance gain carries a well-documented cost, though: such models are widely regarded as ``black boxes,'' and clinical adoption depends on being able to inspect, not merely trust, the reasoning behind a prediction \cite{ref8}. Surveys by Tjoa and Guan \cite{ref31}, van der Velden et~al. \cite{ref10}. And the edited volume of Samek et~al. \cite{ref32} converge on the same caution. Qualitative heatmap visualisation is the most common form of ``explanation'' reported in clinical deep learning. It is routinely treated as sufficient evidence of interpretability on its own, without further testing \cite{ref8}, \cite{ref9}.

Selvaraju et~al. \cite{ref11} proposed Grad-CAM. It weights a convolutional layer's feature maps by the gradient of the target class score. Chattopadhay et~al. \cite{ref12} introduced Grad-CAM++. They replaced globally average-pooled weights with pixel-wise weighted second-order derivatives, giving sharper localisation of small or multiple lesions. Vanilla gradient saliency \cite{ref33} takes the gradient of the class score directly with respect to input pixels. The result is noisier but at pixel resolution. Unlike the CAM family, it needs no target-layer choice. It still requires differentiable, white-box access to gradients, so it is architecture-agnostic, not model-agnostic in the stricter sense that applies to input/output-only methods below. Muhammad and Yeasin \cite{ref34} proposed the class-agnostic, gradient-free Eigen-CAM. It projects a layer's activations onto their dominant principal component. Wang et~al. \cite{ref35} proposed the gradient-free Score-CAM. It weights each activation channel by its measured effect on the model's own output when used as an input mask. This avoids gradient saturation but costs extra compute. Beyond these five, model-agnostic methods such as SHAP \cite{ref36} and LIME \cite{ref37} and gradient-based variants such as Integrated Gradients \cite{ref38} and LayerCAM \cite{ref39} broaden the attribution landscape. But they fall outside the scope of this paper. Section~\ref{sec:methods}.3 gives the full formulation of each of the five methods used here.

A growing body of work argues that visual plausibility alone is not enough for medical XAI. Adebayo et~al. \cite{ref13} found that several popular saliency methods do not respond to model or label randomisation. That sanity check should make any trustworthy method fail loudly. Hooker et~al. \cite{ref14} created the ROAR benchmark, which compares methods by feature-removal faithfulness. Ghassemi et~al. \cite{ref9} and Rudin \cite{ref8} make a related point: post-hoc explanation methods get over-trusted in healthcare without this kind of quantitative scrutiny. Geirhos et~al. \cite{ref40} gave shortcut learning its name and formal definition. Zech et~al. \cite{ref41} showed a real case of it in medical imaging, where a pneumonia classifier partly learned hospital-specific image artefacts instead of pathology. That is exactly the failure mode the shortcut audit in Section~\ref{sec:results}.4 is built to catch. Guo et~al. \cite{ref42} showed that modern deep networks are often miscalibrated even when accuracy is high, which is why Section~\ref{sec:results}.4 includes a calibration analysis. This paper does not just cite this tradition; it puts it into practice for gallbladder ultrasound, running the Adebayo sanity check, a deletion/insertion faithfulness test in the spirit of ROAR, a shortcut audit (all in Section~\ref{sec:results}.4), and a cross-replicate consistency score $O^c$.

Three gaps stand out in the literature above. First, no gallbladder-ultrasound study found in this search quantitatively compares Grad-CAM, Grad-CAM++, Saliency Maps, Eigen-CAM, and Score-CAM against one another on complementary reliability dimensions for this task. RadFormer~\cite{ref19} comes closest, offering quantitative attention-map localization against radiologist annotations, but for a single explanation mechanism rather than a multi-method comparison. Second, none of this literature tests whether its explanations are \textit{reproducible} across independently trained models, as opposed to being an artifact of one training run. Third, the quantitative-XAI-reliability tradition described above---sanity checks, faithfulness testing, shortcut audits, calibration---has not, as far as the available evidence shows, been brought together for gallbladder ultrasound specifically. This paper addresses all three gaps (Table~\ref{tbl:9}, Section~\ref{sec:discussion}.3). It does not claim to be the first application of XAI, or of quantitative XAI evaluation, to gallbladder ultrasound.
